

\documentclass[10pt,twoside,slovak,a4paper]{article}

\usepackage[slovak]{babel}

\usepackage[IL2]{fontenc} 
\usepackage[utf8]{inputenc}
\usepackage{graphicx}
\usepackage{url} 
\usepackage{hyperref} 

\usepackage{cite}
\graphicspath{ {./} }


\pagestyle{headings}

\title{Deep learning, neurónové siete a ich využitie v praxi\thanks{Semestrálny projekt v predmete Metódy inžinierskej práce, ak. rok 2021/22, vedenie: Ing. Zuzana Špitálová}} % 

\author{Marián Gomola\\[2pt]
	{\small Slovenská technická univerzita v Bratislave}\\
	{\small Fakulta informatiky a informačných technológií}\\
	{\small \texttt{xgomola@stuba.sk}}
	}

\date{\small 18. október 2021} 

\begin{document}

\maketitle


\begin{abstract}
V rámci zadania témy "modelovanie v softvérovom inžinierstve," som sa rozhodol pozrieť na špecifický druh strojového učenia a to "deep learning". 
chcem sa sústrediť najmä na neurónové siete ako základný model/ tzv.: “chrbtovú kosť“ deep learningu. 
Opíšem ich vznik, vývoj a vysvetlím ako fungujú. Vysvetlím aj inšpiráciu v ľudských neurónoch.
Uvediem a vysvetlím algoritmus ako príklad. Algoritmus bude v programovacom jazyku Python.
Pozriem sa aj na využitie deep learningu v praxi, ako aj jeho možné využitie v budúcnosti. 
\end{abstract}
\pagebreak
\section{Úvod}


\end{document}